% source: David Mumford, "Abelian Varieties", TIFR Studies in Mathematics 5, Oxford UP, 1970
% pdf_page: 0058
% doc_page: 47 (Chapter II, Section 5, "Cohomology and base change")
% retrieved: 2026-07-11
% transcribed_by: claude-fable-5 (vision, rendered at 200dpi; scanned book, no text layer)

\DeclareMathOperator{\Spec}{Spec}

% [running head: ALGEBRAIC THEORY VIA VARIETIES, page 47]

% [continuation of the proof of the main theorem from the previous page]
Moreover, for all $A$-algebras $B$, $\{U_i \times_Y \Spec B\}$ is an affine
covering of $X \times_Y \Spec(B)$, and
$C^p(\mathfrak{A}, \mathcal{F}) \otimes_A B$ is the module of \v{C}ech
$p$-cochains of $\mathcal{F} \otimes_A B$ for this covering. Therefore
\[
  H^p(X \times_Y \Spec B, \mathcal{F} \otimes_A B)
    \cong H^p(C^{\cdot} \otimes_A B)
\]
for all $B$, and, in fact, functorially in $B$.

We need the following basic lemma

\paragraph{Lemma 1.}
\textit{Let $C^{\cdot}$ be a complex of $A$-modules ($A$ any noetherian ring)
such that the $H^i(C^{\cdot})$ are finitely generated $A$-modules and such
that $C^p \neq (0)$ only if $0 \leq p \leq n$. Then there exists a complex
$K^{\cdot}$ of finitely generated $A$-modules such that $K^p \neq (0)$ only
if $0 \leq p \leq n$ and $K^p$ is free if $1 \leq p \leq n$ and a
homomorphism of complexes $\phi : K^{\cdot} \to C^{\cdot}$ such that $\phi$
induces isomorphisms
$H^i(K^{\cdot}) \xrightarrow{\ \sim\ } H^i(C^{\cdot})$, all $i$. Moreover if
the $C^p$ are $A$-flat, then $K^0$ will be $A$-flat too.}

\paragraph{Proof.}
We define, by descending induction on $m$, diagrams:
\[
\begin{array}{ccccccccc}
 & & K^{m} & \xrightarrow{\ \partial^{m}\ } & K^{m+1}
   & \xrightarrow{\ \partial^{m+1}\ } & K^{m+2} & \longrightarrow & \cdots \\[4pt]
 & & \big\downarrow{\scriptstyle \phi_{m}} & & \big\downarrow{\scriptstyle \phi_{m+1}}
   & & \big\downarrow{\scriptstyle \phi_{m+2}} & & \\[4pt]
\cdots \longrightarrow & C^{m-1} \longrightarrow & C^{m}
   & \xrightarrow[\ \partial^{m}\ ]{} & C^{m+1}
   & \xrightarrow[\ \partial^{m+1}\ ]{} & C^{m+2} & \longrightarrow & \cdots
\end{array}
\]
Put $K^p = 0$ for $p > n$. Suppose we have defined
$(K^p, \phi_p, \partial^p)$ for $p \geq m + 1$ such that the following
conditions hold:
\begin{itemize}
  \item[(i)] $\partial^p \phi_p = \phi_{p+1} \partial^p$, $(p \geq m + 1)$.
  \item[(ii)] $\partial^{p+1} \circ \partial^p = 0$, $(p \geq m + 1)$.
  \item[(iii)] The $\phi^p$ induces isomorphisms in cohomology
    $H^q(K^{\cdot}) \xrightarrow{\ \sim\ } H^q(C^{\cdot})$ for
    $q \geq m + 2$, and a surjection
    $\ker \partial^{m+1} \to H^{m+1}(C^{\cdot})$.
  \item[(iv)] The $K^p$ are $A$-free and finitely generated,
    $(p \geq m + 1)$.
\end{itemize}
We then construct $K^m$, $\partial^m$ and $\phi_m$ so as to satisfy
(i)-(iii) with $m + 1$ replaced by $m$.

Suppose first that $m \geq 0$. Let $B^{m+1}$ be the kernel of the
homomorphism $\ker \partial^{m+1} \to H^{m+1}(C^{\cdot})$. Since $B^{m+1}$
is finitely generated
% [proof continues on the next page]
